% -----------------------------------------------------------------------------
% Document class
% scrartcl = clean article class for reports; bibliography in TOC, good tables
% -----------------------------------------------------------------------------
\documentclass[
  bibliography=totoc,
  captions=tableheading,
]{scrartcl}

% Fix minor KOMA compatibility issues
\usepackage{scrhack}

% Warn if recompilation needed (refs, TOC, etc.)
\usepackage[aux]{rerunfilecheck}

% Core math packages (align, symbols, extensions)
\usepackage{amsmath}
\usepackage{amssymb}
\usepackage{mathtools}

% Modern font handling (requires lualatex/xelatex)
\usepackage{fontspec}
\recalctypearea{} % recompute layout after font setup

% Language settings (hyphenation, captions)
\usepackage[english]{babel}

% Modern math fonts + ISO style (clean scientific notation)
\usepackage[
  math-style=ISO,
  bold-style=ISO,
  sans-style=italic,
  nabla=upright,
  partial=upright,
  mathrm=sym,
]{unicode-math}

\setmathfont{Latin Modern Math} % default math font
\setmathfont{XITS Math}[range={scr, bfscr}] % nicer script letters
\setmathfont{XITS Math}[range={cal, bfcal}, StylisticSet=1]

% Numbers & units formatting (\num, \SI, etc.)
\usepackage[
  locale=US,
  separate-uncertainty=true,
  per-mode=symbol-or-fraction,
  detect-all,
]{siunitx}

% Quotes (recommended with biblatex)
\usepackage[english=american]{csquotes}

% Float placement control (figures/tables positioning)
\usepackage{float}
\floatplacement{figure}{htbp}
\floatplacement{table}{htbp}

% Keep figures inside sections (prevents drifting)
\usepackage[section,below]{placeins}

% text wrapping around figures
\usepackage{wrapfig}      

% Caption styling (bold labels, smaller font)
\usepackage[
  labelfont=bf,
  font=small,
  width=0.9\textwidth,
]{caption}

% Subfigures (compare plots side-by-side)
\usepackage{subcaption}

% Include images
\usepackage{graphicx}

% Better tables (clean formatting + number alignment)
\usepackage{tabularray}
\UseTblrLibrary{booktabs, siunitx}

% Typography improvements (spacing, kerning)
\usepackage{microtype}

% Bibliography system (modern replacement for BibTeX)
\usepackage[backend=biber]{biblatex}
\addbibresource{../../../latex/references.bib}

% Clickable links + PDF metadata
\usepackage[
  unicode,
  pdfusetitle,
  pdfcreator={},
  pdfproducer={},
]{hyperref}

% Better PDF bookmarks
\usepackage{bookmark}

% Custom commands (needed for definitions below)
\usepackage{expl3}
\usepackage{xparse}

% Upright e (useful for exp, softmax, etc.)
\ExplSyntaxOn
\NewDocumentCommand \E {} {\symup{e}}
\ExplSyntaxOff

% Upright differential d (clean integrals)
\ExplSyntaxOn
\NewDocumentCommand \dif {m} {\mathinner{\symup{d} #1}}
\ExplSyntaxOff

% Vector command (bold vectors in math mode)
\ExplSyntaxOn
\let\vaccent=\v
\RenewDocumentCommand \v {}{
  \TextOrMath{\vaccent}{\symbf}
}
\ExplSyntaxOff

% Better spacing for scalable brackets
\usepackage{mleftright}

% Shortcuts for scalable brackets: \l( ... \r)
\ExplSyntaxOn
\let\ltext=\l
\RenewDocumentCommand \l {}{\TextOrMath{\ltext}{\mleft}}
\let\raccent=\r
\RenewDocumentCommand \r {}{\TextOrMath{\raccent}{\mright}}
\ExplSyntaxOff

% Absolute value and norm (auto-scaling with *)
\DeclarePairedDelimiter{\abs}{\lvert}{\rvert}
\DeclarePairedDelimiter{\norm}{\lVert}{\rVert}

% Definition equality symbols
\newcommand{\defeq}{\vcentcolon=}
\newcommand{\eqdef}{=\vcentcolon}

% No paragraph indent (clean report style)
\setlength{\parindent}{0pt}

\title{Task 04: Graph Neural Network for 2D Neutrino Interaction Position Reconstruction}
\author{Akram Aki}
\date{\today}

\begin{document}

\maketitle

\section{Introduction}

The goal of this task was to reconstruct the two-dimensional neutrino interaction position from measured Cherenkov light. 
Each event consists of a variable-length set of detected photon hits, where each hit contains the detection time and detector position, and the target 
is the true interaction coordinate $(x_{\mathrm{true}}, y_{\mathrm{true}})$.

\section{Method}

The input hit features and target coordinates were standardized using the mean and standard deviation computed on the training set. 
The same normalization constants were applied to the validation and test sets, and final predictions were transformed back to the 
original coordinate scale in $\mathrm{m}$.

The reconstruction error was evaluated using the Euclidean distance
\begin{equation*}
    d_{\mathrm{err}} = \sqrt{ (x_{\mathrm{pred}} - x_{\mathrm{true}})^2 + (y_{\mathrm{pred}} - y_{\mathrm{true}})^2 }.
\end{equation*}

Each event was represented as a graph whose nodes are photon hits with normalized time and detector position features. 
A Dynamic Edge Convolution GNN with two DynamicEdgeConv layers, global mean pooling, and a final MLP was used to predict the normalized interaction coordinates.

The model was trained with mean squared error loss and the Adam optimizer for 30 epochs. The best checkpoint was selected using the validation loss, 
with the selected checkpoint at epoch 15 and early stopping patience set to 6.

\section{Results}

\autoref{fig:loss-curves} shows the training and validation loss curves. The validation loss decreased during 
training, and the best checkpoint was selected at epoch 15. The test loss shown in the plot is in normalized coordinate space, since the training objective 
was computed on normalized labels.

\begin{figure}[htbp]
    \centering
    \includegraphics[width=0.9\textwidth]{../build/loss_curves.pdf}
    \caption{Training and validation loss curves. The dashed horizontal line indicates the test loss, and the vertical dotted line indicates the selected best 
    checkpoint epoch.}
    \label{fig:loss-curves}
\end{figure}

\autoref{fig:pred-vs-true} compares the predicted and true interaction coordinates on the test set after transforming 
the coordinates back to the original scale. The predictions follow the diagonal trend, indicating that the model learned the main structure of 
the reconstruction task.

\begin{figure}[htbp]
    \centering
    \begin{subfigure}{0.48\textwidth}
        \centering
        \includegraphics[width=\textwidth]{../build/pred_vs_true_x.pdf}
        \caption{$x$-coordinate.}
        \label{fig:pred-vs-true-x}
    \end{subfigure}
    \hfill
    \begin{subfigure}{0.48\textwidth}
        \centering
        \includegraphics[width=\textwidth]{../build/pred_vs_true_y.pdf}
        \caption{$y$-coordinate.}
        \label{fig:pred-vs-true-y}
    \end{subfigure}
    \caption{Predicted and true interaction coordinates on the test set after transforming the labels back to the original coordinate scale in 
    $\mathrm{m}$. The diagonal line indicates perfect reconstruction.}
    \label{fig:pred-vs-true}
\end{figure}

The distribution of reconstruction errors is shown in \autoref{fig:error-hist}. Most events have 
moderate errors, while the tail of the distribution corresponds to events with larger position reconstruction error.

\begin{figure}[htbp]
    \centering
    \includegraphics[width=0.85\textwidth]{../build/position_error_histogram.pdf}
    \caption{Distribution of the Euclidean reconstruction error $d_{\mathrm{err}}$ on the test set in $\mathrm{m}$. The legend reports the mean, 
    standard deviation, SEM, and median. Mean is the average reconstruction error, standard deviation describes the spread of event-wise errors, 
    SEM is the standard error of the mean, and median is the central value of the error distribution.}
    \label{fig:error-hist}
\end{figure}

\newpage

\autoref{fig:error-2d} visualizes 300 randomly selected test events. True and predicted positions are connected by 
lines whose color represents \(d_{\mathrm{err}}\), while the line width also varies slightly with the error. The detector positions are shown as well and were 
determined by collecting the unique detector hit positions from the data. This plot suggests that reconstruction quality may depend on spatial position, 
since errors appeared larger farther away from the detector region.

\begin{figure}[htbp]
    \centering
    \includegraphics[width=0.85\textwidth]{../build/reconstruction_error_2d_colored.pdf}
    \caption{Two-dimensional reconstruction visualization for 300 randomly selected test events. 
    True and predicted positions are connected by lines colored according to the Euclidean reconstruction error 
    $d_{\mathrm{err}}$ in $\mathrm{m}$, as shown by the colorbar. The line width also varies slightly with the error, and detector positions are shown.}
    \label{fig:error-2d}
\end{figure}

\newpage

The spatial dependence of the reconstruction quality is studied more directly in \autoref{fig:error-heatmap}. 
Test events were binned according to their true interaction position $(x_{\mathrm{true}}, y_{\mathrm{true}})$. For each bin, the color represents the 
mean Euclidean reconstruction error, while the number inside the bin gives the standard deviation of the error values in that bin. 
The heatmap confirms that the reconstruction quality varies with spatial position, with larger errors in some outer regions.

\begin{figure}[htbp]
    \centering
    \includegraphics[width=0.85\textwidth]{../build/position_error_heatmap_mean_with_std.pdf}
    \caption{Spatial dependence of reconstruction quality on the test set. Events are binned by true interaction position, 
    the color gives the mean Euclidean reconstruction error in $\mathrm{m}$, and the number inside each bin gives the standard deviation of 
    the error values in that bin. Detector positions are marked in the plot.}
    \label{fig:error-heatmap}
\end{figure}

\section{Conclusion}

The Dynamic Edge Convolution GNN was able to learn the reconstruction task and produced reasonable predictions on the test set. 
The spatial plots show that the reconstruction performance is not uniform over the detector plane, with larger errors in some outer regions.

As possible extensions, the GNN could be compared to baseline approaches such as a simple fully connected network using fixed summary features, 
a point-cloud model, or a model based on handcrafted event-level quantities. Future work could also include fixing random seeds and documenting the exact 
training configuration more carefully for reproducibility. All code used for preprocessing, training, evaluation, and plotting is available in the accompanying 
\href{https://github.com/AkramAki/Advanced_DeepLearning_Seminar}{GitHub repository}.

\printbibliography

\end{document}
